\subsection*{Ex.~3.2 (10 pts)}

No. The hierarchy would survive, because the premise of the question fails: a
device's access time is set by its material \emph{and} its size, and at terabyte
scale it is size that dominates.

Every read must physically cross the array: the address propagates through the
decoder to a word line, and the data returns along a bit line through a sense
amplifier. That time is bounded below by the time needed to cross the array, whose
linear dimension grows as $\sqrt{C}$ in the capacity $C$. Take the fastest cell
actually manufacturable, a high-density 6T SRAM bit at $0.021\,\mu\mathrm{m}^2$: one
terabyte is $8\times10^{12}$ cells covering $1.68\times10^{5}\,\mathrm{mm}^2$, about
2.4 whole 300\,mm wafers and some $0.41\,\mathrm{m}$ on a side. Crossing that and
returning takes at least $2.7\,\mathrm{ns}$ \emph{in vacuum}, eleven clock cycles at
4\,GHz, before a single transistor switches; real on-chip wires are several times
slower than $c$. A terabyte of the fastest material is thus nowhere near as fast as
a kilobyte of it, and no choice of material moves that floor.

The penalty is already visible within one technology. L1, L2 and L3 on a modern CPU
are all SRAM, on the same die at the same process node, yet they answer in roughly
4, 14 and 40 cycles. Only capacity differs: longer word and bit lines carry more RC
delay, deeper decoders add stages, and larger bit-line capacitance makes sense
amplifiers resolve more slowly. So the hierarchy is not an artifact of mixing cheap
slow memory with expensive fast memory. It reappears with the material held fixed.

Power says the same thing. The fastest cells are the leakiest, six transistors per
bit with static current in each and no idle state as cheap as a DRAM capacitor's. A
terabyte of them could not be powered or cooled as one block, and spreading it out
to cool it pays the propagation delay back again.

The right use of such a device would therefore be to put a small fast array in front
of it to absorb the references that locality makes predictable: to build the
hierarchy inside it. Cost per bit would stop being the motive; the device's own
size--speed trade-off would supply one.
